% ══════════════════════════════════════════════════════════════════
% PREAMBULE — Tous les packages du rapport (règles du guide ENASTIC)
% Police Times New Roman 12pt · Marges 2cm · Interligne 1,5 · Justifié
% ══════════════════════════════════════════════════════════════════

% ── Encodage & langue ─────────────────────────────────────────────
\usepackage[utf8]{inputenc}
\usepackage[T1]{fontenc}
\usepackage[french]{babel}
\usepackage{csquotes}

% ── Police Times New Roman (règle 2 du guide) ────────────────────
\usepackage{mathptmx}          % Times New Roman texte + maths

% ── Marges : 2cm partout (règle 2 du guide) ──────────────────────
\usepackage[top=2cm,bottom=2cm,left=2cm,right=2cm]{geometry}
\setlength{\headheight}{15pt}

% ── Interligne 1,5 (règle 2 du guide) ────────────────────────────
\usepackage{setspace}
\onehalfspacing

% ── Paragraphes justifiés (défaut LaTeX) + retrait ───────────────
\usepackage{microtype}
\setlength{\parindent}{1cm}

% ── Couleurs ──────────────────────────────────────────────────────
\usepackage{xcolor}
\usepackage[most]{tcolorbox}
\usepackage[dvipsnames,table]{xcolor}
\definecolor{bleuENASTIC}{RGB}{55,90,170}   % bleu bordure page garde
\definecolor{grisbleu}{RGB}{70,90,110}
% Couleurs pastel des diagrammes UML
\definecolor{vertdiag}{RGB}{225,240,225}
\definecolor{bleudiag}{RGB}{222,235,250}
\definecolor{jaunediag}{RGB}{253,246,220}
\definecolor{rougediag}{RGB}{250,228,225}
\definecolor{violetdiag}{RGB}{238,230,248}

% ── Graphiques ────────────────────────────────────────────────────
\usepackage{graphicx}
\usepackage{float}
\usepackage{tikz}
\usetikzlibrary{arrows.meta,shapes.geometric,positioning,fit,backgrounds,calc}

% ── Tableaux ──────────────────────────────────────────────────────
\usepackage{booktabs}
\usepackage{tabularx}
\usepackage{multirow}
\usepackage{longtable}
\usepackage{array}

% ── Encadrés ──────────────────────────────────────────────────────
\usepackage{mdframed}

% ── Listes ────────────────────────────────────────────────────────
\usepackage{enumitem}

% ── Mathématiques ─────────────────────────────────────────────────
\usepackage{amsmath,amssymb}

% ── Code source ───────────────────────────────────────────────────
\usepackage{listings}
\definecolor{codebg}{rgb}{0.97,0.97,0.97}
\definecolor{codeblue}{rgb}{0,0,0.8}
\definecolor{codegreen}{rgb}{0,0.5,0}
\definecolor{codepurple}{rgb}{0.5,0,0.5}
\definecolor{codegray}{rgb}{0.5,0.5,0.5}
\lstset{
	backgroundcolor=\color{codebg},
	basicstyle=\ttfamily\small,
	keywordstyle=\color{codeblue}\bfseries,
	commentstyle=\color{codegreen}\itshape,
	stringstyle=\color{codepurple},
	numberstyle=\tiny\color{codegray},
	numbers=left,numbersep=8pt,
	frame=single,breaklines=true,
	tabsize=4,showstringspaces=false,
	captionpos=b,language=Python
}

% ── Abréviations ──────────────────────────────────────────────────
\usepackage{acronym}

% ── Hyperliens ────────────────────────────────────────────────────
% Liens cliquables : citations [n] en bleu -> bibliographie,
% URL/DOI en bleu -> consultation directe en ligne,
% ET numeros de page du Sommaire/listes cliquables (linktoc=all)
\definecolor{bleuLien}{RGB}{0,60,160}
\usepackage[colorlinks=true,
linkcolor=black,        % renvois internes (sections, figures)
citecolor=bleuLien,     % citations [1] cliquables et visibles
urlcolor=bleuLien,      % URL et DOI cliquables
linktoc=all,            % Sommaire/listes : titre ET numero de page cliquables
bookmarksnumbered=true, % signets PDF numerotes (1, 1.1, 1.2...)
bookmarksopen=true]{hyperref}

% ── En-têtes / pieds de page ──────────────────────────────────────
\usepackage{fancyhdr}
\usepackage{emptypage}
\pagestyle{fancy}
\fancyhf{}
\fancyhead[L]{\small\leftmark}
\fancyhead[R]{\small PFE Licence --- 2025/2026}
\fancyfoot[C]{\thepage}
\renewcommand{\headrulewidth}{0.4pt}

% ── Bibliographie IEEE ────────────────────────────────────────────
% hyperref=true : citations [n] du corps cliquables -> bibliographie ;
% backref=false ; les champs url/doi de chaque entrée sont cliquables -> site
\usepackage[backend=biber,style=ieee,sorting=none,
hyperref=true,defernumbers=false]{biblatex}
\addbibresource{references.bib}

% ── Légendes : figure DESSOUS, tableau DESSUS (règle 3) ──────────
\usepackage{caption}
\captionsetup[figure]{labelfont=bf,labelsep=colon,position=below}
\captionsetup[table]{labelfont=bf,labelsep=colon,position=above}

% ── Numérotation Figure 2.4 / Tableau 1.5 (règle 3) ──────────────
\numberwithin{figure}{chapter}
\numberwithin{table}{chapter}
\numberwithin{equation}{chapter}

% ── CHAPITRES EN MAJUSCULES, SEULS AU CENTRE D'UNE PAGE (règle 1) ─
\usepackage{titlesec}
% Le titre du chapitre : taille 14, MAJUSCULES, centré sur page vierge
\titleformat{\chapter}[display]
{\normalfont\centering}
{\fontsize{14}{18}\selectfont\bfseries\MakeUppercase{\chaptertitlename}~\thechapter}
{12pt}
{\fontsize{14}{18}\selectfont\bfseries\MakeUppercase}
\titlespacing*{\chapter}{0pt}{0.32\textheight}{20pt}

% Commande : titre du chapitre SEUL au centre d'une page vierge
% (règle 1 du guide) — le contenu commence à la page suivante.
% Usage : \pagechapitre{chap:label}{Titre du chapitre}
\newcommand{\pagechapitre}[2]{%
	\clearpage
	\chapter{#2}\label{#1}%
	\thispagestyle{empty}%
	\vspace*{\fill}%
	\clearpage
}